\documentclass[11pt,reqno]{amsart}
\usepackage{amsmath,amssymb,amsthm}
\usepackage[margin=1in]{geometry}
\usepackage{booktabs}
\usepackage[colorlinks=true,linkcolor=blue,citecolor=blue,urlcolor=blue]{hyperref}

\theoremstyle{plain}
\newtheorem{theorem}{Theorem}
\newtheorem{proposition}[theorem]{Proposition}
\newtheorem{lemma}[theorem]{Lemma}
\newtheorem{corollary}[theorem]{Corollary}
\theoremstyle{definition}
\newtheorem{definition}[theorem]{Definition}
\theoremstyle{remark}
\newtheorem{remark}[theorem]{Remark}

\newcommand{\Q}{\mathbb{Q}}
\newcommand{\R}{\mathbb{R}}
\newcommand{\C}{\mathbb{C}}
\newcommand{\Z}{\mathbb{Z}}
\newcommand{\F}{\mathbb{F}}
\newcommand{\Sph}{\mathbb{S}}
\newcommand{\Bool}{\mathbb{B}}
\newcommand{\m}{\mathfrak{m}}
\newcommand{\Oo}{\mathcal{O}}
\newcommand{\dg}{^{\dagger}}
\newcommand{\Gal}{\operatorname{Gal}}
\newcommand{\spec}{\operatorname{spec}}
\newcommand{\Spec}{\operatorname{Spec}}
\newcommand{\Tr}{\operatorname{Tr}}
\newcommand{\Char}{\operatorname{char}}
\DeclareMathOperator{\GL}{GL}

\title[A characteristic-zero obstruction to a geometric proof of RH]{Two no-go theorems for the geometric strategy\\ on the Riemann Hypothesis}
\author{James Scully}
\address{}
\email{jamesjscully@gmail.com}
\date{June 2026}

\begin{document}

\begin{abstract}
We study the obstruction to importing Weil's function-field proof of the Riemann Hypothesis (RH) into the
number-field setting. We reduce RH to a single \emph{reality condition} $F\dg=\bar F$ on a polarized
arithmetic Hodge category, and prove two no-go results. The \emph{Geometric No-Go} (Theorem~\ref{thm:B}) is
a structural theorem: a similitude $F\cdot F\dg=[\deg]$ \emph{underdetermines the signature} of its invariant
Hermitian form, so the positive-definiteness that is RH is independent of the duality structure and cannot be
forced by it. The \emph{Foundational No-Go} (Theorem~\ref{thm:A}) shows that a finite additive characteristic
at the archimedean place of $\Q$ is \emph{contradictory} with $\Q$ being a field of characteristic $0$;
consequently the finite ground that Weil's proof requires at $\infty$ is not merely missing but forbidden.
The unifying mechanism is a \emph{count/radius decoupling}: the reality is multiplicative and free at every
place, while the scale $q=|\F_q|$ is additive (a characteristic) and exists only over a genuine finite field.
We neither prove nor disprove RH; we establish that if RH is provable, the proof cannot route through giving
the archimedean place a finite ground.
\end{abstract}

\maketitle

\section{Introduction}\label{sec:intro}

Over a smooth projective curve $C/\F_q$, RH for the zeta function $\zeta_C$ is \emph{Weil's theorem}
\cite{Weil}: the Frobenius eigenvalues on $H^1$ satisfy $|\alpha|=q^{1/2}$ (purity), forced by the
\emph{Hodge index theorem} on the surface $C\times_{\F_q}C$---a geometric positivity that holds
\emph{because} the surface is projective over the finite field $\F_q$. The ``geometric strategy'' for RH over
$\Q$ is to construct the analogous projective surface over a finite ground for $\Spec\Z$; this is the program
of the field with one element / absolute geometry \cite{Soule,Deitmar,Borger,ConnesConsaniF1}. Our purpose is
to prove that this strategy is foreclosed.

Informally: in the function-field world \emph{every} place sits over a finite field of constants. The places
of $\Q$ also sit over finite fields---the residue field at a prime $p$ is $\F_p$---with one exception, the
archimedean place, whose completion is $\R$. The real numbers are infinite and have no finite ground. Weil's
proof needs a finite ground at \emph{every} place; $\Q$ has one everywhere but at infinity. We show this gap
is irreparable: it is the definitional identity \emph{archimedean $\equiv$ characteristic $0$}, and a finite
ground there would force a prime number to equal $0$.

We prove two independent results.

\begin{itemize}
\item \textbf{Theorem~\ref{thm:A} (Foundational No-Go).} A finite additive characteristic at the archimedean
place $\infty$ of $\Q$ is contradictory with $\Char\Q=0$. The finite ground at $\infty$ the strategy requires
cannot exist.
\item \textbf{Theorem~\ref{thm:B} (Geometric No-Go).} A similitude $F\cdot F\dg=[\deg]$ underdetermines the
signature of its invariant Hermitian form; positive-definiteness (which is RH for the datum) is an
independent condition that no structure seeing only the duality data can force.
\end{itemize}

\noindent\emph{Calibration.} Theorem~\ref{thm:A} is a theorem resting on Ostrowski's theorem, the definition
of characteristic, and the finite-semifield classification. Theorem~\ref{thm:B} is proved as a structural
statement (Theorem~\ref{thm:B}); that \emph{every named coherence framework} reduces to its hypotheses is a
case-by-case verification (\S\ref{sec:exhaustion}), not a single meta-theorem. We do not claim to prove or
disprove RH.

\section{The reduction to a reality condition}\label{sec:reduction}

We summarize the reduction; the construction is carried out and validated numerically against the
function-field case elsewhere. One builds a polarized category with a representing object $\Sph_\Z$ carrying:
a Frobenius $F$ with $\Tr(F^t\mid H^1)=\sum_\rho e^{t\rho}$ (zeros $=\spec F$); a $(p,q)$-grading at $\infty$;
a structural duality $D$ realizing the functional equation; and a \emph{categorical} polarization $P$ equal
to the canonical pairing of $D$ (not a chosen metric). On the function-field anchor this reproduces Weil's
proof, with the categorical polarization equal to the Gram pairing. Peeling the resulting single positivity
to its atom yields the following equivalence.

\begin{proposition}\label{prop:reality}
RH for the datum is the \emph{reality condition}
\[
F\dg=\bar F,
\]
the arithmetic (Rosati) adjoint equals the analytic (complex-conjugation) adjoint; equivalently
$|\alpha|^2=\deg$ for every eigenvalue. The three classical formulations---the reality condition,
Hodge--Riemann metric-compatibility, and Hilbert--P\'olya self-adjointness---coincide on this condition.
\end{proposition}

\section{The count/radius decoupling}\label{sec:decoupling}

The condition $F\dg=\bar F$ factors, and the factorization is the mechanism behind both no-gos.

\begin{itemize}
\item \textbf{$F\cdot F\dg=[\deg]$ (the functional equation) is free}: a formal consequence of
rigidity/duality, holding at every radius (including off-line). It is the \emph{multiplicative} datum---the
involution, $\Gal$, the reality.
\item \textbf{$|\alpha|=\sqrt q$ (the radius) is the deep half}: $q/\alpha=\bar\alpha$. It is the
\emph{additive} datum, $q=|\F_q|$ being the order of the additive group, i.e.\ the characteristic.
\end{itemize}

\noindent Over $\F_q$ the two are married by a single Frobenius $x\mapsto x^q$ acting on $+$ and $\times$ at
once, with $q$ the one number tying additive cardinality to multiplicative degree. At $\infty$ the marriage
breaks: one obtains the reality without the scale. This is the invariant that recurs below.

\section{The Geometric No-Go}\label{sec:B}

\begin{definition}
A \emph{Frobenius datum} is a finite-dimensional complex space $V$ with a non-degenerate Hermitian form
$H$ and an endomorphism $F\in\GL(V)$ that is a \emph{similitude of factor $c>0$}: $F^*HF=cH$, equivalently
$F\cdot F\dg=c\cdot 1$ for the adjoint $\dg$ of $H$. RH for the datum is $|\alpha|=\sqrt c$ for every
eigenvalue $\alpha$ of $F$ (\emph{Step B}); the similitude/functional equation is \emph{Step A}. The spectrum
of a similitude is stable under the involution $\sigma\colon\alpha\mapsto c/\bar\alpha$.
\end{definition}

\begin{theorem}[Geometric No-Go]\label{thm:B}
Let $F$ be a similitude of factor $c>0$ on $V$, with $\sigma$-stable spectrum. Then:
\begin{enumerate}
\item\emph{(Step A is free.)} A non-degenerate $F$-invariant Hermitian form exists for \emph{every}
$\sigma$-stable spectrum, in particular for spectra with $|\alpha|\neq\sqrt c$.
\item\emph{(Step B $\Leftrightarrow$ RH.)} Such a form can be chosen positive-definite if and only if every
$|\alpha|=\sqrt c$. If some $|\alpha|\neq\sqrt c$, every $F$-invariant non-degenerate Hermitian form is
indefinite, with negative index equal to the number of $\sigma$-pairs $\{\alpha,c/\bar\alpha\}$ off the
circle.
\item\emph{(No-go.)} Consequently the similitude structure does not determine the signature;
positive-definiteness (Step B) is an independent condition, and any structure that constructs $F$ from its
duality data alone forces only Step A. Forcing Step B requires imposing positive-definiteness as an extra
axiom (a $C^*$-structure, a positive fibre functor), which is logically equivalent to assuming RH for the
datum.
\end{enumerate}
\end{theorem}

\begin{proof}
If $Fx=\alpha x$ then $c\,H(x,x)=H(Fx,Fx)=|\alpha|^2H(x,x)$, so $(|\alpha|^2-c)H(x,x)=0$: either
$|\alpha|^2=c$ or $x$ is isotropic ($H(x,x)=0$). If $H$ is positive-definite no nonzero vector is isotropic,
forcing $|\alpha|^2=c$ for all eigenvalues; conversely if some $|\alpha|^2\neq c$ the eigenvector is
isotropic, so $H$ is not definite. This is~(2). For existence and the signature: diagonalize $F$; the
involution $\sigma$ pairs the spectrum. On a $\sigma$-fixed eigenvalue ($|\alpha|^2=c$) set $H$ definite
($+1$) on the eigenline; on a $\sigma$-pair $\{\alpha,\beta=c/\bar\alpha\}$ with $\alpha\neq\beta$ the two
eigenvectors are isotropic and the unique invariant pairing couples them, a hyperbolic plane of signature
$(1,1)$. Summing yields a non-degenerate invariant $H$ for any $\sigma$-stable spectrum, with signature
$(d,0)\oplus(1,1)^{\oplus k}$, $d=\#$ on-circle eigenvalues and $k=\#$ off-circle pairs; this proves~(1) and
completes~(2). Part~(3) is immediate.
\end{proof}

\begin{remark}[Numerical witness]
For $c=7$, an on-line datum yields $F^*HF=cH$ to $1.8\times10^{-15}$ with signature $(4,0,0)$
(positive-definite); an off-line datum with two pairs off the circle yields the \emph{same} similitude
residual $1.8\times10^{-15}$ with signature $(4,2,0)$ (indefinite). Step A is free at both radii; Step B holds
only on-line.
\end{remark}

\subsection{Exhaustion: every coherence structure is an instance}\label{sec:exhaustion}
Each coherence structure proposed for RH presents the similitude/Step-A data and is silent on the signature,
so off-line data are consistent in it and Step B is an external input. Table~\ref{tab:exhaustion} records, for
each, what is free and why Step B is not forced.

\begin{table}[h]
\small
\begin{tabular}{@{}lll@{}}
\toprule
structure & Step A (free) & why Step B is not forced \\
\midrule
Tannakian / rigid tensor & finite-dim, duality, FE & off-line $F$ generates a consistent \\
 & & non-polarizable category \cite{DeligneMilne} \\
$C^*$ / Hermitian & invariant form, non-degenerate & signature $(d,0)$ on-line vs.\ $(k,k)$ Krein; \\
 & & Doplicher--Roberts \emph{consumes} positivity \cite{DR} \\
o-minimal / tame & proper-pushforward purity over $\F_q$ & $\Gamma_\R$ has infinite-discrete poles, \\
 & & hence non-o-minimal \cite{BBT,vdD} \\
Weil / Li positivity & the explicit-formula form & PSD $\Leftrightarrow$ RH (circular) \cite{Bombieri} \\
Bost--Connes QSM & KMS states & free energy is RH-silent \cite{BostConnes} \\
sym$^2$ / Rankin--Selberg & Euler-borne, independent & forces the edge $\mathrm{Re}\,s=1$, not the line \\
de Bruijn--Newman flow & dynamics + zero statistics & the GUE lever is universal/RH-silent \\
prismatic / $q$-de Rham & $q$-deformed FE & finite-prime-indexed, archimedean-blind \\
condensed / six-functor & coherent duality & supplies Step A only \\
abelian-variety / Rosati & $F\cdot F\dg=[\deg]$ tautology & the scalar is free; $F\dg=\bar F$ is RH \\
\bottomrule
\end{tabular}
\medskip
\caption{Each row exhibits the free Step~A and the off-line consistency predicted by Theorem~\ref{thm:B}.}
\label{tab:exhaustion}
\end{table}

\section{The Foundational No-Go}\label{sec:A}

We use only standard facts. For a place $v$ of $\Q$: the completion $\Q_v$; the residue field
$\kappa(v)=\Oo_v/\m_v$ when $v$ is non-archimedean ($\Oo_v$ the valuation ring); and $\Char$ the additive
order of $1$. The Weil hypothesis at a place is ``$\kappa(v)$ is a finite field $\F_{q_v}$ and Frobenius acts
with eigenvalue-modulus $q_v^{w/2}$.''

\begin{lemma}[Ostrowski {\cite{Ostrowski}}]\label{lem:ostrowski}
Every nontrivial place of $\Q$ is equivalent to exactly one of: the $p$-adic absolute value $|\cdot|_p$ (one
per prime $p$; non-archimedean; $\Q_p$; $\kappa=\F_p$), or the usual absolute value $|\cdot|_\infty$
(archimedean; $\Q_\infty=\R$).
\end{lemma}

\begin{lemma}\label{lem:equiv}
For a nontrivial place $v$ of $\Q$ the following are equivalent: \emph{(i)} $v$ is non-archimedean;
\emph{(ii)} $\Z$ is bounded in $\Q_v$; \emph{(iii)} $v$ has a finite residue field $\kappa(v)$;
\emph{(iv)} $\Char\kappa(v)=p>0$. Dually, $v$ is archimedean iff the Archimedean property holds
($\forall x\,\exists n: |n\cdot 1|_v>|x|_v$) iff $\Q_v\supseteq\R\supseteq\Q$ iff $\Char(\Q_v)=0$.
\end{lemma}

\begin{theorem}[Foundational No-Go]\label{thm:A}
There is no object furnishing the archimedean place $\infty$ of $\Q$ with a finite additive characteristic
(equivalently, with a finite residue field). The Weil hypothesis fails at $\infty$, and not contingently: it
is contradictory with $\Char\Q=0$.
\end{theorem}

\begin{proof}
By Lemma~\ref{lem:ostrowski}, $\infty$ is archimedean and $\Q_\infty=\R\supseteq\Q$. Suppose a finite
additive characteristic $p>0$ were furnished at $\infty$, i.e.\ $p\cdot 1=0$ in the residue/completion datum
$k$ carrying the place, where $\Q\subseteq k$ (every completion of $\Q$ contains $\Q$; for $\infty$,
$\Q\subseteq\R=k$). Restricting to $\Q\subseteq k$ gives $p\cdot 1=0$ in $\Q$, contradicting $\Char\Q=0$.
Equivalently, by Lemma~\ref{lem:equiv} a finite additive characteristic at $\infty$ is equivalent to $\infty$
being non-archimedean, contradicting its classification as \emph{the} archimedean place.
\end{proof}

\begin{proposition}[Characteristic-1 diagnostic]\label{prop:char1}
Drop the field axiom to admit a finite additive characteristic on a semifield. The only finite semifield with
idempotent addition (``characteristic $1$'') is the Boolean semifield $\Bool=(\{0,1\},\max,\cdot)$, whose
multiplicative group has order $1$ \cite{ConnesConsaniChar1,Lescot}. Hence any finite-additive-characteristic
structure at $\infty$ avoiding a finite field is characteristic-$1$, forcing the scale $q=1$. More directly: a
finite scale $q\geq2$ has $q-1\geq1$ units forming the cyclic group $\mu_{q-1}$ of $(q-1)$-th roots of unity,
whose coexistence with a finite-characteristic additive structure is exactly a finite field $\F_q$ of
characteristic $p$. Thus a finite additive characteristic and a nontrivial finite multiplicative scale coexist
only on a genuine finite field.
\end{proposition}

\begin{corollary}\label{cor:A}
Weil's proof of RH over $\F_q(C)$ requires $C$ projective over the finite field $\F_q$, hence a finite residue
field at every place of $\F_q(C)$ (including its place at infinity, a finite extension of $\F_q$). The
number-field analogue requires a finite residue field at every place of $\Q$; by Theorem~\ref{thm:A} this
fails at $\infty$, and by Proposition~\ref{prop:char1} it is irreparable by weakening ``field'' to
``semifield.'' The Weil/Hodge-index strategy does not transfer to $\Q$.
\end{corollary}

\begin{remark}[The identity]
Lemmas~\ref{lem:ostrowski}--\ref{lem:equiv} and Theorem~\ref{thm:A} amount to the chain of definitional
equivalences
\[
\text{archimedean}\ \equiv\ \kappa\supseteq\Q\ \equiv\ \Char 0\ \equiv\ n\cdot 1\not\to 0\ \equiv\
\text{all primes inverted}\ \equiv\ \text{non-ultrametric},
\]
and a finite additive scale at $\infty$ negates the last four clauses, hence the first.
\end{remark}

\section{The \texorpdfstring{$\F_1$}{F1}-exhaustion}\label{sec:f1}

Theorem~\ref{thm:A} is corroborated by an exhaustive survey of candidate finite grounds, each failing by the
count/radius decoupling. Among bases: the sphere spectrum $\Sph$ (homotopy cardinality is a count, not the
radius \cite{ConnesConsaniRR}); tropical / characteristic-$1$ (dequantizes, $q\to1$); Borger--Witt
(finite-prime-indexed \cite{Borger}); monoid schemes \cite{Deitmar,Soule}; the Tate-twist period ($2\pi$, a
transcendental). Among degenerations: the archimedean Fargues--Fontaine / twistor curve \cite{FarguesFontaine}
has residue field $\C$, $\Gal(\C/\R)$ of modulus $1$, slope $\tfrac12$---the reality, no scale; the $p$-adic
Fargues--Fontaine curve carries $q$ only by importing $\Spec\F_p$ via Witt vectors, and $\R$ has no finite
residue field to import.

\section{Scope, consequences, status}\label{sec:scope}

\noindent\textbf{Proven.} Theorem~\ref{thm:A} and Proposition~\ref{prop:char1} are theorems. The reduction
(Proposition~\ref{prop:reality}) and the count/radius decoupling are established constructively and validated
on the function-field anchor.

\noindent\textbf{Theorem-shaped.} Theorem~\ref{thm:B} is proved (\S\ref{sec:B}); that each named coherence
framework instantiates it is verified case by case in Table~\ref{tab:exhaustion}, not by one meta-theorem.

\noindent\textbf{Not claimed.} We neither prove nor disprove RH. Theorem~\ref{thm:A} is not an undecidability
statement; it is the precise limit of one strategy. That the characteristic divide obstructs RH is the
founding intuition of the $\F_1$ program; the contribution here is its formalization as a contradiction and
the exhaustion showing every geometric route hits it.

\noindent\textbf{Consequence.} If RH is provable, the proof cannot route through a finite ground at $\infty$
(closed by Theorem~\ref{thm:A}). RH must be approached on the archimedean/characteristic-$0$ side---analytic
number theory, the explicit formula, de Branges / Hilbert--P\'olya \emph{as analysis} (a self-adjoint
operator built without an underlying finite-ground geometry)---or by a foundation reconceiving the identity
\emph{archimedean $\equiv$ characteristic $0$}, which no current framework achieves.

\subsection*{Acknowledgments}
The results were obtained through an extended computational research program: a disciplined,
falsification-driven exploration that constructed and tested candidate objects against the function-field
anchor and a kill-test suite, with numerical witnesses for the constructive claims. The author thanks that
process for catching several over-claims en route to the statements above.

\begin{thebibliography}{99}

\bibitem{BBT} B.~Bakker, Y.~Brunebarbe, J.~Tsimerman, \emph{o-minimal GAGA and a conjecture of Griffiths},
Invent.\ Math.\ (2022).

\bibitem{Bombieri} E.~Bombieri, \emph{Remarks on Weil's quadratic functional in the theory of prime numbers},
Rend.\ Mat.\ Acc.\ Lincei (2000); and the explicit-formula positivity $W(f)\geq0\Leftrightarrow$ RH (Weil,
1952).

\bibitem{Borger} J.~Borger, \emph{$\Lambda$-rings and the field with one element}, arXiv:0906.3146.

\bibitem{BostConnes} J.-B.~Bost, A.~Connes, \emph{Hecke algebras, type III factors and phase transitions with
spontaneous symmetry breaking in number theory}, Selecta Math.\ (1995).

\bibitem{ConnesConsaniChar1} A.~Connes, C.~Consani, \emph{Characteristic one, entropy and the absolute point},
arXiv:0911.3537.

\bibitem{ConnesConsaniF1} A.~Connes, C.~Consani, \emph{Schemes over $\F_1$ and zeta functions},
Compos.\ Math.\ (2010).

\bibitem{ConnesConsaniRR} A.~Connes, C.~Consani, \emph{Riemann--Roch for $\overline{\Spec\Z}$ over the sphere
spectrum} (2023).

\bibitem{ConreyLi} J.~B.~Conrey, X.-J.~Li, \emph{A note on some positivity conditions related to zeta and
$L$-functions}, Int.\ Math.\ Res.\ Not.\ (2000).

\bibitem{Deitmar} A.~Deitmar, \emph{Schemes over $\F_1$}, in \emph{Number Fields and Function Fields}
(2005).

\bibitem{DeligneMilne} P.~Deligne, J.~Milne, \emph{Tannakian categories}, in \emph{Hodge Cycles, Motives, and
Shimura Varieties}, LNM 900 (1982).

\bibitem{DR} S.~Doplicher, J.~Roberts, \emph{A new duality theory for compact groups}, Invent.\ Math.\ 98
(1989).

\bibitem{FarguesFontaine} L.~Fargues, J.-M.~Fontaine, \emph{Courbes et fibr\'es vectoriels en th\'eorie de
Hodge $p$-adique}, Ast\'erisque 406 (2018).

\bibitem{Lescot} P.~Lescot, \emph{Alg\`ebre absolue}, Ann.\ Sci.\ Math.\ Qu\'ebec (2009).

\bibitem{Ostrowski} A.~Ostrowski, \emph{\"Uber einige L\"osungen der Funktionalgleichung
$\varphi(x)\varphi(y)=\varphi(xy)$}, Acta Math.\ 41 (1916). See also J.~Neukirch, \emph{Algebraic Number
Theory}, Springer (1999), II.\S3.

\bibitem{Soule} C.~Soul\'e, \emph{Les vari\'et\'es sur le corps \`a un \'el\'ement}, Mosc.\ Math.\ J.\ (2004).

\bibitem{vdD} L.~van den Dries, \emph{Tame Topology and o-minimal Structures}, LMS Lecture Notes 248 (1998).

\bibitem{Weil} A.~Weil, \emph{Sur les courbes alg\'ebriques et les vari\'et\'es qui s'en d\'eduisent},
Hermann (1948).

\end{thebibliography}

\end{document}
